\bookchapter{第五十六章：打磨与发布——测试、调试与上架 Google Play}

\chapterintro{App 做出来了，然后呢？这一章回答“从能跑”到“能交付”的三件事：怎么高效地调试找 bug、怎么写测试防回归、怎么把 App 送上 Google Play 让全世界下载。上架不只是仪式感——它是把一个作品变成产品的分水岭。}

\booksection{一、调试三板斧：Logcat、断点、UI 检查器}

\begin{center}
\begin{tabularx}{\linewidth}{|X|X|X|}
\hline
\textbf{工具} & \textbf{干什么} & \textbf{怎么用} \tabularnewline
\hline
Logcat & 看系统/自己打的日志 & 代码里 Log.d("TAG","值=\$x")，面板里搜 TAG 过滤 \tabularnewline
断点调试 & 让程序停在某行，看变量值 & 行号旁点一下出红点，Debug 模式运行，逐步走 \tabularnewline
Layout Inspector & 看运行中的界面结构 & 菜单 Tools Layout Inspector，实时看组件树 \tabularnewline
\hline
\end{tabularx}
\end{center}

\begin{center}\small\color{gray}调试三板斧\end{center}

\begin{tcbinfo}{小贴士：报错怎么读？}
崩溃后 Logcat 里会有红色堆栈（Stacktrace），从下往上找\textbf{第一个自己包名开头的行}（比如 com.example.yourapp），那行基本就是案发现场。先看“Exception 类型”（空指针？类型转换？），再看“at ... 你的文件:行号”。这招能解决 80\% 的崩溃。

\end{tcbinfo}

\booksection{二、测试：给代码上保险}

手测靠不住——改了这个功能，另一个悄悄坏了。测试就是“后悔药”兼“防弹衣”。安卓测试分三层，新手先吃透第一层：

\begin{center}
\begin{tabularx}{\linewidth}{|X|X|X|X|}
\hline
\textbf{层级} & \textbf{测什么} & \textbf{放哪里} & \textbf{速度} \tabularnewline
\hline
单元测试 & 纯逻辑（计算、校验） & app/src/test/ & 秒级 \tabularnewline
仪器测试 & 真实设备上的界面行为 & app/src/androidTest/ & 分钟级 \tabularnewline
手动测试 & 真人感受、体验 & —— & 永远需要 \tabularnewline
\hline
\end{tabularx}
\end{center}

\begin{center}\small\color{gray}测试三层\end{center}

\begin{lstlisting}[language=Kotlin]
// 把"金额字符串转数字"这种逻辑抽成纯函数
fun parseAmount(input: String): Double? {
    val v = input.trim().toDoubleOrNull()
    return if (v != null && v > 0) v else null
}

// 测试它（JUnit4）
@Test
fun emptyInputReturnsNull() {
    assertEquals(null, parseAmount(""))
}
@Test
fun negativeReturnsNull() {
    assertEquals(null, parseAmount("-5"))
}
@Test
fun normalValueParses() {
    assertEquals(12.5, parseAmount("12.5"))
}
\end{lstlisting}

记账本里的金额校验、统计计算，都该有单元测试罩着。跑测试：右键 test 目录 Run。绿色全过，放心改代码。

\booksection{三、签名打包：给 App 盖上“印章”}

安卓要求所有安装包必须\textbf{签名}——就像盖公章，证明“这个包是我出的，没被篡改”。调试时 Android Studio 自动用 debug 签名；发布要用自己的 release 签名。流程：

\begin{enumerate}
  \item 菜单 Build Generate Signed Bundle / APK；
  \item 创建 keystore 文件（密钥库）——\textbf{设置密码并妥善保管}；
  \item 选择 release + AAB 格式，构建完成。
\end{enumerate}

\textbf{避坑提醒：keystore 丢了 = App 丢了}\par

上架后更新 App 必须用\textbf{同一个}签名，否则商店会拒绝更新。keystore 文件和密码一旦丢失，你只能换个新应用重新上架——老用户全都接不上。务必把 keystore 和密码备份到安全的地方（云盘加密 + 记事本），这是安卓开发者最昂贵的教训之一。

\booksection{四、上架 Google Play 全流程}

注册开发者账号

准备商店素材

上传 AAB

填数据安全表

审核（数小时\textasciitilde{}数天）

上架

\begin{center}\small\color{gray}图 6：Google Play 上架六步走\end{center}

\begin{itemize}
  \item \textbf{注册}：一次性 25 美元注册费，需能用的 Google 账号；
  \item \textbf{素材}：应用图标（512×512）、截图（至少 2 张）、功能描述、隐私政策链接（收集数据必填）；
  \item \textbf{数据安全表}：如实申报收集了哪些数据——不联网的记账本可以"零收集"；
  \item \textbf{审核}：新账号首次上架常被人工审核，数小时到数天不等，被打回就按意见改。
\end{itemize}

国内用户注意：Google Play 需要特殊网络环境，且中国大陆地区分发通常走国内商店（华为、小米、应用宝等）。\textbf{上架国内商店的流程类似，但各商店有自己的审核口味}——同一个 AAB/APK 可以多店分发。作为练手，先走通一家就赢了。

\booksection{五、商店素材怎么做：第一印象就是下载量}

商店页是用户“逛店”时看到的橱窗。三个素材决定第一印象：

\begin{itemize}
  \item \textbf{图标}：512×512，别用截图凑合。没设计能力就用 AI 工具生成一个简洁几何图形，\textbf{在手机桌面上缩小后还能一眼认出}才算合格；
  \item \textbf{截图}：至少 2 张，最好 4-6 张，展示最核心的功能场景；每张配一句人话文案（"一秒记一笔"比"高效的记账解决方案"好十倍）；
  \item \textbf{描述}：第一句说清"这是什么 + 给谁用"，别堆砌"最专业、最强"——商店搜索权重主要看关键词和下载转化，真诚比夸张有效。
\end{itemize}

\begin{tcbinfo}{小贴士：A/B 看看同行}
上架前逛一圈商店里同类 App 的前十名，截图对比：他们图标什么风格？截图讲什么卖点？标题怎么写？\textbf{模仿不是抄袭}——学套路、换自己的内容，这是所有产品冷启动最快的一课。

\end{tcbinfo}

\booksection{五、性能优化与崩溃上报小清单}

App 上架前，这两件事值得各花半小时：

\booksubsection{性能自查四条}

\begin{enumerate}
  \item \textbf{列表用 LazyColumn}：长列表千万别用 Column 硬排——懒加载是安卓性能的第一课；
  \item \textbf{图片别硬塞原图}：用 Coil 图片库自动缩放缓存，一张 5MB 原图足以卡爆低端机；
  \item \textbf{避免在主线程干重活}：数据库、网络、大文件都在协程里做；
  \item \textbf{用 Profiler 量一量}：Android Studio 自带的性能分析器能看内存、CPU、耗电——先测量，再优化，别凭感觉。
\end{enumerate}

\booksubsection{崩溃上报}

你测试时没崩，不代表用户手机不崩。上架前接入\textbf{Firebase Crashlytics}（谷歌官方，免费）：几行配置，用户崩溃时自动上报堆栈，你在后台就能看到“哪个机型、哪个版本、崩在哪一行”。\textbf{没有崩溃监控的上架，等于闭着眼睛开车}。

\begin{tcbpractice}{来动手：把记账本"交付"一次}
(1) 给 parseAmount 等逻辑补 3 个单元测试；(2) 用 Logcat 亲手抓一个自己制造的崩溃并修好；(3) 生成一个 release 签名的 APK，发给朋友安装（这才是“作品”的仪式感）；(4) 有余力就注册开发者账号，把记账本真正提交上架。前三条今晚就能做完。

\end{tcbpractice}

\begin{tcbsummary}{本章小结}
\begin{itemize}
  \item 调试三板斧：Logcat 看日志、断点看变量、Layout Inspector 看界面
  \item 读崩溃：从堆栈下往上找自己包名的第一行
  \item 测试三层：单元（秒级，先学这个）/仪器/手动
  \item 发布必须签名：keystore 和密码备份好，丢了 App 就"换人"了
  \item 上架六步：注册 素材 上传 AAB 数据安全表 审核 上架
\end{itemize}

\end{tcbsummary}
